\section{Proof of Theorem~\ref{thm:sgem}}
\label{sec:proof}

This appendix collects the lemmas and the full proof of
Theorem~\ref{thm:sgem}. We follow the notation of
Section~\ref{subsec:theory}: $\theta = (\boldsymbol{\alpha}_{1:K}, \mathbf{R})
\in \Theta = \mathcal{A}^K \times (\Delta_K)^M$ with $\mathcal{A}$
finite, $\varepsilon_t$ is the per-iteration surrogate error
(in the loss scale; see~\eqref{eq:eps_t}), $N = \sum_m |\mathcal{C}_m|$,
and the variational free energy
parameterised by a loss function $u$ is
\begin{equation*}
\mathcal{F}_u(q, \theta) \;=\; \sum_{m=1}^{M} |\mathcal{C}_m|
\sum_{k=1}^{K} q_{mk}\, \log \frac{r_{mk}\, e^{-u(\boldsymbol{\alpha}_k, \mathbf{r}_k)}}{q_{mk}}.
\end{equation*}
We write $L_u(\theta) = \max_q \mathcal{F}_u(q, \theta)$ for the
size-weighted log-likelihood induced by loss $u$, and abbreviate
$L_{\mathrm{true}} = L_{u_{\mathrm{true}}}$,
$L^{(t)} = L_{u^{(t)}}$.

\subsection{Preliminary Lemmas}

\begin{proposition}[ELBO maximisation]
\label{prop:elbo}
For every loss function $u: \mathcal{A} \times [0, 1]^M \to \mathbb{R}_{\geq 0}$
and every $\theta\in\Theta$,
$L_u(\theta) = \max_{q\in\mathcal{Q}} \mathcal{F}_u(q, \theta)$ with
the maximum attained at
\begin{equation*}
q^{\star}_{mk}(\theta) \;=\; \frac{r_{mk}\,e^{-u(\boldsymbol{\alpha}_k, \mathbf{r}_k)}}
{\sum_{j} r_{mj}\,e^{-u(\boldsymbol{\alpha}_j, \mathbf{r}_j)}}.
\end{equation*}
\end{proposition}

\begin{proof}
Fix $\theta$ and $m$. The factor $|\mathcal{C}_m|$ is a positive
constant; it can be pulled out of the inner maximisation in $q_{m,:}$
without affecting the optimum. Writing $a_{mk} = r_{mk}\,
e^{-u(\boldsymbol{\alpha}_k, \mathbf{r}_k)}$ and using
$\sum_k q_{mk} = 1$, the inner expression equals
$-\mathrm{KL}(q_{m,:} \| a_{m,:}/\!\sum_j a_{mj}) + \log\sum_j a_{mj}$.
The KL-divergence is non-negative and vanishes exactly at
$q_{mk} = a_{mk}/\sum_j a_{mj}$, which yields the claim. Summing
over $m$ with weights $|\mathcal{C}_m|$ recovers
$L_u(\theta) = \sum_m |\mathcal{C}_m|\,\log\sum_k a_{mk}$.
\end{proof}

A consequence is the Neal--Hinton decomposition
$L_u(\theta) = \mathcal{F}_u(q, \theta) + \sum_m |\mathcal{C}_m|\,
\mathrm{KL}\!\bigl(q_{m,:}\|q^{\star}_{m,:}(\theta)\bigr)$,
implying $L_u(\theta) \geq \mathcal{F}_u(q,\theta)$ with equality iff
$q = q^{\star}(\theta)$.

\begin{lemma}[Uniform ELBO deviation]
\label{lem:elbo-dev}
Under Assumption~\ref{ass:bounded}, for every $q$, $\theta$, $t$,
$\bigl|\mathcal{F}_{u_{\mathrm{true}}}(q,\theta) -
\mathcal{F}_{u^{(t)}}(q,\theta)\bigr| \leq N\,\varepsilon_t$.
\end{lemma}

\begin{proof}
The terms $\log r_{mk}$ and $-\log q_{mk}$ cancel between the two
ELBOs, leaving
\begin{equation*}
\mathcal{F}_{u_{\mathrm{true}}} - \mathcal{F}_{u^{(t)}}
\;=\; \sum_m |\mathcal{C}_m| \sum_k q_{mk}\,
\bigl[u^{(t)}(\boldsymbol{\alpha}_k, \mathbf{r}_k) -
u_{\mathrm{true}}(\boldsymbol{\alpha}_k, \mathbf{r}_k)\bigr].
\end{equation*}
Bounding the inner bracket by $\varepsilon_t$ and using
$\sum_k q_{mk} = 1$ gives
$|\mathcal{F}_{u_{\mathrm{true}}} - \mathcal{F}_{u^{(t)}}|
\leq \varepsilon_t \sum_m |\mathcal{C}_m| = N\,\varepsilon_t$.
\end{proof}

\begin{corollary}
\label{cor:lik-dev}
For every $\theta$, $|L_{\mathrm{true}}(\theta) - L^{(t)}(\theta)|
\leq N\,\varepsilon_t$.
\end{corollary}

\begin{proof}
Apply Lemma~\ref{lem:elbo-dev} to the maxima of the two ELBOs and
use $|\max_q f - \max_q g| \leq \sup_q |f - g|$.
\end{proof}

Finally, we record three auxiliary bounds --- on the softmax map, on
the responsibilities, and on the Q-function --- used to analyse the
limit points in Part~(b) of the main proof; they compare the
surrogate-based and the true responsibilities.

\begin{lemma}[Softmax stability]
\label{lem:softmax}
Let $a_k, \tilde a_k > 0$ satisfy
$e^{-\varepsilon} a_k \leq \tilde a_k \leq e^{\varepsilon} a_k$. Then
the normalised vectors $p_k = a_k/\sum_j a_j$ and
$\tilde p_k = \tilde a_k/\sum_j \tilde a_j$ obey
$|\tilde p_k - p_k| \leq (e^{2\varepsilon} - 1)\,p_k \leq 3\varepsilon\,p_k$
for $\varepsilon \leq 1$.
\end{lemma}

\begin{proof}
$\tilde p_k \leq a_k e^{\varepsilon}/\sum_j a_j e^{-\varepsilon}
= p_k\,e^{2\varepsilon}$ and similarly $\tilde p_k \geq p_k\,e^{-2\varepsilon}$.
\end{proof}

\begin{lemma}[Responsibility deviation]
\label{lem:resp}
Let $w_{mk}(\theta) = q^{\star}_{mk}(\theta)$ be the
true-surrogate responsibilities (Proposition~\ref{prop:elbo} with
$u = u_{\mathrm{true}}$) and $q^{(t)}_{mk}$ the surrogate
responsibilities of the E-step at iteration $t$. For
$\varepsilon_t \leq 1$,
$|q^{(t)}_{mk} - w_{mk}(\theta^{(t)})| \leq 3\varepsilon_t\,
w_{mk}(\theta^{(t)})$.
\end{lemma}

\begin{proof}
Apply Lemma~\ref{lem:softmax} with
$a_{mk} = r^{(t)}_{mk}\,e^{-u_{\mathrm{true}}(\boldsymbol{\alpha}^{(t)}_k, \mathbf{r}^{(t)}_k)}$
and $\tilde a_{mk}$ obtained by replacing $u_{\mathrm{true}}$ by
$u^{(t)}$. The ratio $\tilde a_{mk}/a_{mk}
= e^{u_{\mathrm{true}} - u^{(t)}}$ satisfies
$e^{-\varepsilon_t} \leq \tilde a_{mk}/a_{mk} \leq e^{\varepsilon_t}$
by~\eqref{eq:eps_t}, so Lemma~\ref{lem:softmax} applies with
$\varepsilon = \varepsilon_t$.
\end{proof}

\begin{lemma}[Q-function deviation]
\label{lem:qdev}
Define the surrogate Q-function
\begin{equation*}
\tilde Q^{(t)}(\theta;\theta') \;=\; \sum_{m=1}^M |\mathcal{C}_m|
\sum_{k=1}^K q^{(t)}_{mk}(\theta')
\bigl[\log r_{mk} \,-\, u^{(t)}(\boldsymbol{\alpha}_k, \mathbf{r}_k)\bigr].
\end{equation*}
Under Assumption~\ref{ass:bounded} there exists a constant
$C_Q = C_Q(N, K, U_{\max})$ such that for every $\theta, \theta'$ and
every $t$ with $\varepsilon_t \leq 1$,
$|\tilde Q^{(t)}(\theta;\theta') - Q_{\mathrm{true}}(\theta;\theta')|
\leq C_Q\,\varepsilon_t$.
\end{lemma}

\begin{proof}
Add and subtract
$\sum_m |\mathcal{C}_m| \sum_k q^{(t)}_{mk}\bigl[\log r_{mk}
- u_{\mathrm{true}}(\boldsymbol{\alpha}_k, \mathbf{r}_k)\bigr]$.
The first piece, the surrogate-error term, is
$\sum_m |\mathcal{C}_m| \sum_k q^{(t)}_{mk}
(u_{\mathrm{true}} - u^{(t)})$ and is bounded by $N\varepsilon_t$.
The remainder splits as
$\sum_m |\mathcal{C}_m| \sum_k (q^{(t)}_{mk} - w_{mk})(-u_{\mathrm{true}})$
and
$\sum_m |\mathcal{C}_m| \sum_k (q^{(t)}_{mk} - w_{mk})\log r_{mk}$.
The first is bounded by $3NK\,U_{\max}\,\varepsilon_t$ via
Lemma~\ref{lem:resp} (because $|u_{\mathrm{true}}| \leq U_{\max}$).
For the second, using $|x\log x|\leq 1/e$ on $[0,1]$ together with
the softmax-stability bound of Lemma~\ref{lem:softmax}, one obtains a
bound of $3N\varepsilon_t/e$. Summing the three contributions gives
the claim with $C_Q = N(1 + 3K\,U_{\max} + 3/e)$.
\end{proof}

\subsection{Main Proof}

\begin{proposition}[Quasi-monotonicity of $L_{\mathrm{true}}$]
\label{prop:quasi-mono}
Under Assumptions~\ref{ass:bounded}--\ref{ass:gem}, for every $t$,
\begin{equation}
\label{eq:quasi-mono}
L_{\mathrm{true}}(\theta^{(t+1)}) \;\geq\;
L_{\mathrm{true}}(\theta^{(t)}) \;-\; 2N\,\varepsilon_t.
\end{equation}
\end{proposition}

\begin{proof}
We chain five inequalities. First,
$L_{\mathrm{true}}(\theta^{(t+1)}) \geq
\mathcal{F}_{u_{\mathrm{true}}}(q^{(t)}, \theta^{(t+1)})$ since
$L_u(\theta) = \max_q \mathcal{F}_u(q,\theta)$. Second, by
Lemma~\ref{lem:elbo-dev},
$\mathcal{F}_{u_{\mathrm{true}}}(q^{(t)}, \theta^{(t+1)})
\geq \mathcal{F}_{u^{(t)}}(q^{(t)}, \theta^{(t+1)}) - N\varepsilon_t$.
Third, the GEM Assumption~\ref{ass:gem} yields
$\mathcal{F}_{u^{(t)}}(q^{(t)}, \theta^{(t+1)})
\geq \mathcal{F}_{u^{(t)}}(q^{(t)}, \theta^{(t)})$. Fourth, the
E-step makes the latter equal to $L^{(t)}(\theta^{(t)})$. Fifth,
Corollary~\ref{cor:lik-dev} gives $L^{(t)}(\theta^{(t)}) \geq
L_{\mathrm{true}}(\theta^{(t)}) - N\varepsilon_t$. Combining the
five steps yields~\eqref{eq:quasi-mono}.
\end{proof}

\textbf{Part (a): convergence of the likelihood.}
Set $L_{\max} = \sup_{\theta\in\Theta} L_{\mathrm{true}}(\theta)$,
which is finite because $\Theta$ is compact (Assumption~\ref{ass:cont}
and finiteness of $\mathcal{A}$). Define
$V_t = L_{\max} - L_{\mathrm{true}}(\theta^{(t)}) \geq 0$. From
Proposition~\ref{prop:quasi-mono},
\begin{equation*}
V_{t+1} \;\leq\; V_t \;+\; 2N\varepsilon_t,
\qquad \sum_{t=0}^{\infty} 2N\varepsilon_t < \infty
\;\;\text{by Assumption~\ref{ass:summable}.}
\end{equation*}
The Robbins--Siegmund theorem~\cite{robbins1971convergence} (with
$a_t = 0$, $b_t = 2N\varepsilon_t$) implies that $\{V_t\}$ converges,
hence $\{L_{\mathrm{true}}(\theta^{(t)})\}$ converges. \qed

\textbf{Part (b): stationarity of limit points.}

\textbf{Step 1: existence.} $\Theta$ is compact, so
$\{\theta^{(t)}\}$ has at least one limit point
$\theta^{\star} = \lim_{j\to\infty} \theta^{(t_j)}$.

\textbf{Step 2: contradiction setup.} Suppose $\theta^{\star}$ is not
stationary: there is $\theta'\in\Theta$ and $\delta > 0$ with
$Q_{\mathrm{true}}(\theta';\theta^{\star}) -
Q_{\mathrm{true}}(\theta^{\star};\theta^{\star}) = \delta > 0$.

\textbf{Step 3: continuity.} The map $\theta'\mapsto Q_{\mathrm{true}}(\theta;\theta')$
is continuous on $\Theta$: $w_{mk}(\theta')$ is continuous because
$u_{\mathrm{true}}$ is continuous in $\mathbf{r}$
(Assumption~\ref{ass:cont}) and softmax is continuous; and the
product $w_{mk}\log r_{mk}$ extends continuously to $r_{mk} = 0$
because $w_{mk} = O(r_{mk})$ and $r\mapsto r\log r$ is continuous at
$0$. Thus, for $j$ large enough,
$Q_{\mathrm{true}}(\theta';\theta^{(t_j)}) -
Q_{\mathrm{true}}(\theta^{(t_j)};\theta^{(t_j)}) \geq \delta/2$.

\textbf{Step 4: transfer to $\tilde Q$.} By Lemma~\ref{lem:qdev},
\begin{equation*}
\tilde Q^{(t_j)}(\theta';\theta^{(t_j)}) -
\tilde Q^{(t_j)}(\theta^{(t_j)};\theta^{(t_j)}) \;\geq\;
\delta/2 - 2C_Q\,\varepsilon_{t_j} \;\geq\; \delta/4
\end{equation*}
for $j$ large enough, since $\varepsilon_{t_j}\to 0$.

\textbf{Step 5: GEM ascent.} Because the M-step searches over a
candidate set that contains $\theta'$ (by finiteness of $\mathcal{A}$
and the closure properties of the $\mathbf{R}$-update),
$\tilde Q^{(t_j)}(\theta^{(t_j+1)};\theta^{(t_j)}) \geq
\tilde Q^{(t_j)}(\theta';\theta^{(t_j)})$. Translating back through
Lemma~\ref{lem:qdev} once more,
\begin{equation*}
\Delta_{t_j} := Q_{\mathrm{true}}(\theta^{(t_j+1)};\theta^{(t_j)}) -
Q_{\mathrm{true}}(\theta^{(t_j)};\theta^{(t_j)}) \;\geq\;
\delta/4 - 2C_Q\,\varepsilon_{t_j} \;\geq\; \delta/8
\end{equation*}
for $j$ large enough.

\textbf{Step 6: summability contradiction.} The Gibbs inequality
$L_{\mathrm{true}}(\theta^{(t+1)}) - L_{\mathrm{true}}(\theta^{(t)})
\geq \Delta_t$ together with Proposition~\ref{prop:quasi-mono}
implies $\Delta_t \geq -2N\varepsilon_t$. Hence the non-negative
sequence $\Delta_t + 2N\varepsilon_t$ satisfies
$\sum_t (\Delta_t + 2N\varepsilon_t) < \infty$ by
Robbins--Siegmund applied to $V_t$, and consequently
$\sum_t \Delta_t < \infty$. But Step~5 gives $\Delta_{t_j} \geq
\delta/8$ for infinitely many $j$, yielding $\sum_t \Delta_t = +\infty$
--- a contradiction. Therefore $\theta^{\star}$ is stationary. \qed

\subsection{Convergence Rate}

A by-product of the proof of Theorem~\ref{thm:sgem}(a) is the
explicit rate
\begin{equation}
L_{\mathrm{true}}(\theta^{\star}) - L_{\mathrm{true}}(\theta^{(T)})
\;\leq\; 2N \sum_{t = T}^{\infty} \varepsilon_t.
\end{equation}
If $\varepsilon_t = O(t^{-p})$ with $p > 1$, this gives
$L_{\mathrm{true}}(\theta^{(T)}) \to L_{\mathrm{true}}(\theta^{\star})$
at rate $O(T^{1 - p})$, matching the rate observed for stochastic-EM
analyses~\cite{cappe2009online}.
